\documentclass[a4paper,11pt,DIV=11,BROC=5mm,bigheadings,idxtotoc,cleardoubleempty,halfparskip,oneside,openright]{scrreprt} % 

\usepackage{datetime}

\usepackage{graphicx}
\usepackage{xcolor}

\usepackage{blindtext}

\usepackage{cite}% Zitieren
\usepackage{bibgerm}% Literatur in Deutscher DIN
\usepackage{url}


\usepackage{amsmath}
\usepackage{amssymb}

\usepackage{listings}

\usepackage[utf8]{inputenc}
\usepackage{lmodern}

\usepackage{microtype}


\usepackage{scrpage2} 	% Kopf & Fußzeile im KOMA Stil
\pagestyle{scrheadings}	% Aktiviert Verwendung vordefinierter Kolumnentitel
\clearscrheadfoot 		% alle Standard-Werte und Formatierungen löschen
\setkomafont{pagehead}{\scshape}	% Schriftart in Kopfzeile, \scshape = Kapitelchen
\automark[chapter]{section} % [linke Seite]{rechte Seite}
%\ohead{\def\pagestyle{PDTS}{\hrulefill\includegraphics[width = 6cm]{bilder/thi_logo_quer_cropped}}}
\ohead{\includegraphics[width = 6cm]{bilder/thi_logo_quer_cropped}}

%\ihead{\textsc{Abschlussarbeit}}
\ihead{\headmark}

%\setheadwidth[0pt]{textwithmarginpar}
\ofoot{\vspace{-0.3cm} \pagemark} 						
\ifoot{\vspace{-0.3cm} Dominik Gunther Florian Schlecht} 
				
%\setheadtopline{2pt}				
%\setheadsepline{.4pt}
\setfootsepline{.4pt}	% Trennlinie Fußzeile und Textkörper

%-------------------------------------------------------------






%------------------------------------------------------------------
%% Längenanpassungen
%------------------------------------------------------------------

\setlength{\headsep}{10mm}				% Textabstand zur Kopfzeile
\setlength{\footskip}{15mm}				% Abstand zur Fußzeile
\setlength{\parindent}{0em}				% Einzug nach Absatz

%------------------------------------------------------------------
%% THI - Farbdefinitionen
%------------------------------------------------------------------

% Farben der THI
\definecolor{THIblue}{rgb}{0.0078,0.1176,0.4705}
%%-------------Allgemeine Definitionen----------------------------------
% Farbige Aufwertung der berschriften
\addtokomafont{chapter}{\color{THIblue}}
\addtokomafont{section}{\color{THIblue}}
\addtokomafont{caption}{\color{THIblue}}
\addtokomafont{subsection}{\color{THIblue}}
\addtokomafont{subsubsection}{\color{THIblue}}
\setkomafont{captionlabel}{\color{THIblue}}

%------------------------------------------------------------------
%% Interessante optinale Packages
%------------------------------------------------------------------

% -> Verlinkter Text im PDF
	\usepackage[
			pdftex,
			linkcolor=THIblue,			% Farbe der Verlinkung
			linktocpage=true,			% Im TOC wird Seitenzahl verlinkt(true),bzw. Text(false)
			colorlinks=true,			% 'true' keinen Kasten um Link
	%		citecolor=blue,
	%		pdfhighlight=/P,
	%		bookmarks,
	%		hyperfigures=true,
	%		citebordercolor={0 0 1},
	%		linkbordercolor={0 0 1},
	%		menubordercolor={0 0 1},
	%		backref=true,
	%		pagebackref=true,
	%%		bookmarksopen,
	%		bookmarksnumbered,
	%		pdfpagelabels=false,
	%		pdfstartpage=1,
	%		pdfstartview=Fit,			% Modus beim Öffnen (Fit = An Seitengröße anpassen)
			pdftitle={Umgehen von PID und VID basierten USB-Policies},
			pdfauthor={Dominik Schlecht},
	%		pdfstartview=Fit,
	%		pdfdisplaydoctitle=true,
	%		plainpages=false
				]{hyperref}  



%------------------------------------------------------------------
%% Wichtige Definition für Aufteilung von Formelverzeichnis und Abkürzungsverzeichnis
%------------------------------------------------------------------

% Nomenklaturverzeichnis, Formelzeichenliste Anpassungen für nomenclature: damit lassen sich zwei getrennte Symbolverzeichnisse anlegen, ziehmlich cool!
\usepackage[intoc,compatible,german]{nomencl}	
		\renewcommand{\nomgroup}[1]{	\ifthenelse{\equal{#1}{A}}{\item[{\normalfont\sffamily\bfseries\LARGE\textcolor[rgb]{0,0.112,0.47}{Abkürzungen{\phantom{\Huge $\frac{\frac{\frac{A}{a}}{a}}{\frac{a}{a}}$}}}}]}{	\ifthenelse{\equal{#1}{A}}{\item[{\normalfont\sffamily\bfseries\LARGE\textcolor[rgb]{0,0.112,0.47}{Formelzeichen{\phantom{\Huge $\frac{A}{\frac{a}{a}}$}}}}]}{}}}
		
%% z.B. (1)  \nomenclature{$v_0$}{Startgeschwindigkeit}   -> Formelverzeichnis
\nomenclature[A]{SAF}{Security Awareness Framework}



%------------------------------------------------------------------
%% Anpassung von Abständen, Längen vom Nomenclaturverzeichnis (Abkürzungs- und Formelverzeichnis)
%------------------------------------------------------------------

		% Abstand zwischen Einträgen im Symbolverzeichnis (-\parsep = 0)
		\setlength{\nomitemsep}{-\parsep} 
		\setlength{\nomlabelwidth}{5em}	
		\renewcommand{\nomlabel}[1]{#1 \dotfill}  % Punkte im zwischen Nummer und Kapiteleintrag






%------------------------------------------------------------------
%% Definition Standardordner der Bilddatein für schnellen Zugriff
%------------------------------------------------------------------
%
%\newcommand\figpfad{bilder/}
%


%%-------------Silbentrennung--------------------------------------
\hyphenation{}

%%-------------Index-----------------------------------------------
\makeindex
\makenomenclature

%------------------------------------------------------------------
% Angabe der zu verwendenden Dateien, sodass nur z.B. das aktuelle 
% Dokument compiliert wird; den Rest mit % auskommentieren
		\includeonly{%
			chapter/title_page,%
			chapter/thanks_statement,
            chapter/einleitung,
			chapter/appendix%
		}


%------------------------------------------------------------------------

%-----------------------------------------------------------------
%---------------Dokumentenbeginn----------------------------------
%-----------------------------------------------------------------
	\begin{document}
			
			
% Titelseite
\begin{titlepage}

\phantom{tmpText}

\vspace{1cm}

\begin{figure}[h!]
\centering
\includegraphics[width=\textwidth]{bilder/thi_logo_cropped.pdf}
\end{figure}

  \begin{center}

%\vspace{1cm}
    
    
    \textbf{{\large Seminararbeit/Whitepaper} \\[3ex]
    {\LARGE Umgehen von USB-Deskriptor basierten USB-Richtlinien am Beispiel einer Virtual Desktop Infrastrukture Umgebung} \\[1ex]
    %
    \vfill
    %
    angefertigt von} \\
    \begin{tabular}{ll}
    	Name: & Dominik Gunther Florian Schlecht\\
    	Matrikelnummer: & 00032209\\[2ex]
    	\multicolumn{2}{c}{\textbf{Betreuer:}}\\
    	Technische Hochschule Ingolstadt: & Dr. Paul Stanislaus \\
      Allianz Deutschland AG: & Dr. Ulrich Stremmel und Stephan Gerhager
    \end{tabular}\\[2ex] %Vorname Nachname
    %
    \vfill
    %
    Ingolstadt, \today
  \end{center}
\end{titlepage}
			
				\cleardoublepage

			
%------------------------------------------------------------------------

% Nummerierung
\pagenumbering{roman}

% Sperrvermerk
\chapter*{Sperrvermerk}
Die vorliegende Arbeit \glqq Umgehen von USB-Deskriptor basierten USB-Policies am Beispiel einer virtuellen Umgebung\grqq\ wurde mit Unterstützung der Allianz Deutschland AG erstellt. Es sind jedoch keine internen, vertraulichen oder streng vertraulichen Informationen enthalten. Die Weitergabe des Inhalts der Arbeit im Gesamten oder in Teilen sowie das Anfertigen von Kopien oder Abschriften -- auch in digitaler Form -- sind grunds\"{a}tzlich erlaubt.
\addcontentsline{toc}{chapter}{Sperrvermerk}

% Erkl\"{a}rung
\chapter*{Erkl\"{a}rung}
Hiermit erkl\"{a}re ich, dass ich die vorliegende Seminararbeit bis auf die offizielle Betreuung durch den Aufgabensteller selbstst\"{a}ndig und ohne fremde Hilfe verfasst habe.\par
Die verwendeten Quellen sowie die verwendeten Hilfsmittel sind vollst\"{a}ndig angegeben. W\"{o}rtlich \"{u}bernommene Textteile und \"{u}bernommene Bilder und Zeichnungen sind in jedem Einzelfall kenntlich gemacht. \\[10ex]
Ingolstadt, \today
\addcontentsline{toc}{chapter}{Erkl\"{a}rung}

% Danksagung (optional)
\chapter*{Danksagung}
An dieser Stelle möchte ich mich bei allen Bedanken, die mich bei der Erstellung dieser Arbeit unterstützt haben. Besonderer Dank dabei gebührt den Kollegen aus der Allianz Deutschland AG und meinen Eltern.

\begin{flushright}
\sffamily Dominik Schlecht, \today
\end{flushright}
\addcontentsline{toc}{chapter}{Danksagung}


%------------------------------------------------------------------------

				\cleardoublepage	
					
			% Inhaltsverzeichnis
			%\setcounter{\secnumdepth}{2} % Zwei Ebenen im Inhaltsverzeichnis auflisten
			\tableofcontents
				\cleardoublepage	
		
			% Abkrzungsverzeichnis
			\printnomenclature
				\cleardoublepage	
			
%------------------------------------------------------------------------
			% Neunummerierung des Hauptteils
			\pagenumbering{arabic}
			\setcounter{page}{1}            
            \cleardoublepage
            
            \chapter{Einleitung}
In Zeiten von \textit{Heartbleed}\cite{Heartbleed} und  \textit{Shellshock}\cite{Shellshock}, \textit{Snowden} und der \textit{NSA}\cite{Snowden} und der fortlaufenden Digitalisierung der Industrie und Gesellschaft wird das Thema Informationssicherheit immer wichtiger. Daten werden, unabhängig davon, ob diese Privatpersonen oder Unternehmen zugeordnet sind, immer wertvoller. So ergeben sich beispielsweise aus einem gehackten Smartphone einer Privatperson Informationen wie E-Mail-Adressen, Kontakte und Chat-Verläufe bis hin zu Passwörter für Online-Banking oder persönlichen Bildern. Wenn diese Informationen auf dem Schwarzmarkt verkauft oder online veröffentlicht werden, kann dies für die Personen oft Reputations wie auch finanzielle Schäden nach sich ziehen. Diese Tätigkeiten werden oft unter dem Schlagwort "Cybercrime" zusammengefasst. Betrachtet man Unternehmen, so ist der mögliche finanzielle Schaden wesentlich höher als für Privatpersonen. Durch die Entwendung von Kreditkarten erlitten zum Beispiel mehrere Supermärkte in den USA beträchtliche Reputationsschäden \cite{HackKreditkarten}\cite{HackKreditkarten2}. Eventuell noch höhere Schäden könnte es nach sich ziehen, wenn streng vertrauliche Dokumente von Unternehmen, wie z.B. Konstruktionsskizzen für ein neues Automodell, Quellcode oder vorläufige Geschäftsberichte durch Hacker erbeutet und an ein Konkurrenzunternehmen verkauft würden. Dies hört sich irreal an, aber die Firma McAffee schätzt den Verlust für die Wirtschaft durch "Cybercrime" im Jahr 2014 auf bis zu 575 Milliarden USD\cite{McAffee}. 
Um diesem Trend entgegen zu wirken, müssen Unternehmen Maßnahmen ergreifen, welche das Schutzniveau erhöhen. Oft werden hier auf der technischen Seite nur internetseitige Komponenten beachtet, wie das schnelle Patching von Servern. Dies ist in Hinsicht auf  Poodle und Shellshock sicherlich auch notwendig, jedoch sollte man alle Wege, über welche Daten von Dritten in das Unternehmen gelangen, Daten an Dritte weitergegeben werden könnten, sowie auch interne Bedrohungen wahrnehmen, einschätzen und eindämmen. Eine solche Prüfung war die Grundlage für dieses Dokument.
                \cleardoublepage
                
			\chapter{Motivation}
			
			\chapter{Allgemeine Gefahren bei USB}
			\section{Viren}
			\subsection{Viren auf Treiberebene und neue Möglichkeiten durch Virtualisierungsumgebungen}
			\subsection{Normale Viren}
			\section{Datenabfluss}
			\section{Dateneinschleußung}

			\chapter{Policies}
			\section{Allgemein}
			\section{UBS-Policies}
						
			
			\chapter{Felder bei USB}
			\section{Allgemein}
			\section{PID}
			\section{VID}
			
			\chapter{Umgehung der Policies}
			\section{Teensy}
			\section{Vorgehen}
			\section{Proof of Concept}

			\chapter{Ausblick in die Zukunft}			
			
			\chapter{Gegenmaßnahmen}
	
			
%------------------------------------------------------------------------
			
			% Anhang
			
%------------------------------------------------------------------------
% Anhänge
\appendix
\chapter{Quellcode}

\definecolor{gray}{rgb}{0.5, 0.5, 0.5}
\lstset{% general command to set parameter(s)
	basicstyle=\small, % print whole listing small
	keywordstyle=\color{blue}\bfseries,% underlined boldblack keywords
	identifierstyle=, % nothing happens
	commentstyle=\color{green}, % white comments
	stringstyle=\color{red}\ttfamily, % typewriter type for strings
	showstringspaces=false, % no special string spaces
	numbers=left,
	numberstyle=\color{gray},
	numbersep=5pt,
	captionpos=t,
	breaklines=true}
	
\lstset{language=Perl}

Lorem ..

			
			% Literaturliste
			\newpage
			\bibliographystyle{plaindin}	% DIN Norm für Literaturdarstellung
			\bibliography{ref/ref_liste}	% Pfad und Datei der Ref-Datenbank
		
	\end{document}
%-----------------------------------------------------------------
